% Inhoud van één labo. Gebruikt door de syllabus van dit labo en door de syllabus voor labo 3-5.
\textbf{Doel:} data aggregeren met \code{groupby}, tabellen koppelen met
\code{merge}, herschikken met \code{melt} en \code{pivot}, en eenvoudige
kenmerken van een tijdreeks berekenen.

\begin{itemize}
  \item \textbf{Herhaling: een kapotte cel herstellen} -- Een kleine opkuis-pipeline met twee fouten repareren.
  \item \textbf{\code{groupby} + \code{agg}} -- Statistieken per groep, leeftijdscategorieën met \code{pd.cut}, meerdere benoemde aggregaties tegelijk.
  \item \textbf{\code{merge}} -- Inner en left joins; waarom een dubbele atleet twee keer meetelt, en waarom je dubbels verwijdert vóór het koppelen.
  \item \textbf{\code{melt} en \code{pivot}} -- Lange en brede vorm; vooruitgang tussen twee testmomenten.
  \item \textbf{Tijdreeksen} -- Hartslagzones als \% van $HR_{max}$ (220 $-$ leeftijd), tijd per zone, \code{rolling}, \code{resample}.
  \item \textbf{Minicase per twee} -- Een trainingsrapport: één rij per atleet.
\end{itemize}

\textbf{Oefeningen:} (0) herhaling: de kapotte cel herstellen;
(1) groeperen per geslacht en leeftijdscategorie; (2) wat heeft het
verwijderen van de dubbels gekost?; (3) vooruitgang in spronghoogte;
(4) minuten in zone 4--5 per atleet; \textbf{minicase:}
trainingsrapport met de minuten in zone 4--5, de maximale snelheid en de
afgelegde afstand per atleet.
